\documentclass[10pt,a4paper]{amsart}
\usepackage[margin=19mm]{geometry}
\usepackage{amsmath,amssymb,mathtools}
\usepackage[scr=rsfs,cal=euler]{mathalfa}
\usepackage{xfrac}
\usepackage[unicode,hidelinks,bookmarks=false]{hyperref}
\newcommand{\ie}{i.e.\ }
\newcommand{\cf}{cf.\ }
\newcommand{\resp}{respectively\ }
\newcommand{\Subsection}{\subsection}
\providecommand{\nobreakdash}{\nobreak\hbox{-}}
\setlength{\emergencystretch}{2em}
\multlinegap0pt
\usepackage{bm}
\usepackage{bbm}
\usepackage{dsfont}
\usepackage{xspace}
\usepackage{cancel}
\usepackage{mathtools}
\usepackage{amsthm}
\makeatletter
\@ifundefined{theorem}{\newtheorem{theorem}{Theorem}}{}
\@ifundefined{lemma}{\newtheorem{lemma}[theorem]{Lemma}}{}
\makeatother
\title[Maximal intersecting families revisited]{Maximal intersecting families revisited}
\author{Yongjiang Wu, Yongtao Li, Lihua Feng, Jiuqiang Liu, Guihai Yu}
\date{Source: arXiv:2411.03674v3 (CC BY 4.0)}
\begin{document}
\maketitle

\noindent\emph{Source and licence.} Maximal intersecting families revisited. Version arXiv:2411.03674v3 (CC BY 4.0), distributed under the Creative Commons Attribution 4.0 International licence, as recorded in the arXiv licence metadata for this exact version and shown on the abstract page https://arxiv.org/abs/2411.03674v3. Full rights notice: licenses/arxiv-2411.03674-CC-BY-4.0.txt.\par\smallskip

\noindent\textit{Editorial scope.} Counted unit: the Lemma whose source label is \texttt{Lem52} in the article (source0.tex lines 1343--1351, source label \texttt{Lem52}), delivered together with the article's own complete original proof of it (source0.tex lines 1352--1403, one \texttt{proof} environment of 52 lines). No mathematics is written, repaired or re-derived: statement and proof are the article's own text line for line. Required context retained uncounted: none. Every definition, display and label of those lines is retained unchanged. Omitted from this package and not redistributed: 1--1342, 1404--1559. Nothing inside a retained excerpt is omitted or altered: every retained body is delivered exactly as the article's lines are written, so no omission receipt of any kind accompanies this package. Citations: the retained excerpts carry no citation command at all, so no bibliography entry of the article is transcribed and no reference is redistributed. Reference adaptation: none was needed and none was made; every reference inside the delivered unit resolves inside it, so no unresolved reference or citation is produced. Printed slips kept literal: no printed word, symbol, hypothesis or number of the counted statement or of its proof was altered. Layout and class: the article's own class is \texttt{article} with packages including amsmath, amsfonts, amsthm, amssymb, mathtools, mathrsfs, graphicx, color, latexsym, tikz, calc, tikz-cd, epstopdf, enumerate; the wrapper is the lane's pinned \texttt{amsart} editorial wrapper at 10pt on A4, which restates the theorem-like environments the retained text needs and adds the article's own macro definitions that the retained text needs (none), with extra wrapper packages (bm, bbm, dsfont, xspace, cancel, mathtools). The delivered numbering is the unit's own. Assets: no retained span contains a figure environment, a graphics-inclusion command, a PDF-page inclusion, a table or a TikZ picture; no image file is copied into this package, the package declares no asset dependency and no image file is redistributed with it. Licence: version arXiv:2411.03674v3 of this article is distributed under the Creative Commons Attribution 4.0 International licence, as recorded in the arXiv licence metadata for this exact version and shown on the abstract page https://arxiv.org/abs/2411.03674v3. A full rights notice is delivered at licenses/arxiv-2411.03674-CC-BY-4.0.txt. Characters: every retained excerpt is pure ASCII, so the delivered engine, XeLaTeX with its default Latin Modern text font, reports no missing character.
\begin{lemma}\label{Lem52}
 Let $n \geq 7$ be an integer.  Let $x_1\neq x_2\in [n]$ and $\mathcal{F}, \mathcal{G} \subseteq\left\{A \in\binom{[n]}{3}:  A\cap \{x_1, x_2\}\neq \emptyset\right\}$ be cross-intersecting families. Let $
\mathcal{R}=\left\{A \in\binom{[n]}{3}: \{x_1,x_2\} \subseteq A \right\}$.
Suppose that $\mathcal{R}\subsetneq \mathcal{F}$, $\mathcal{R}\subsetneq \mathcal{G}$ and $(\mathcal{F}\backslash\mathcal{R}) \cap(\mathcal{G}\backslash\mathcal{R}) =\emptyset$.
Then
$$
|\mathcal{F}|+|\mathcal{G}| \leq\binom{n}{3}-\binom{n-3}{3}-\binom{n-4}{2}-\binom{n-5}{1}+2.
$$
\end{lemma}

\begin{proof}
By conditions and symmetry, we may assume that $\{x_1, a, b\}\in \mathcal{F}$ for some
$a\neq b\in [n]\backslash \{x_1, x_2\}$. Then
\begin{align*}
\mathcal{G}\subseteq\left\{G \in\binom{[n]}{3}: x_1\in G\right\}\cup\left\{\{x_2,a, i\}:  i\in [n]\backslash\{x_1,x_2,a\}\right\}\cup\left\{\{x_2,b, i\}:  i\in [n]\backslash\{x_1,x_2,a, b\}\right\}.
\end{align*}
If $\mathcal{F}$ is a star at $x_1$, or $\mathcal{F}$ is a star at $x_1$ and there is at most a $\{x_2, c, d\}\in \mathcal{F}$ for some
$c\neq d\in [n]\backslash \{x_1, x_2\}$, note that $(\mathcal{F}\backslash\mathcal{R}) \cap(\mathcal{G}\backslash\mathcal{R}) =\emptyset$, then we have
$
|\mathcal{F}|+|\mathcal{G}|\leq n-2+1+ \binom{n-1}{2}+n-3+n-4
=\binom{n-1}{2}+3n-8
$.
Note that $\binom{n}{3}-\binom{n-3}{3}-\binom{n-4}{2}-\binom{n-5}{1}+2=\binom{n-1}{2}+\binom{n-2}{2}+3$.
For $n\geq 7$, let
$
g(n)=\binom{n-1}{2}+\binom{n-2}{2}+3-\left(\binom{n-1}{2}+3n-8\right)
=\binom{n-2}{2}-3n+11=\frac{1}{2}\left(n^2-11n+28 \right)
$.
Then $h(n)\geq 0$. So the result follows. By symmetry, the same case holds for $\mathcal{G}$.

Next we consider the following three cases.

(i) If there exist $\{x_2, c, d\}, \{x_2, c, e\}\in \mathcal{F}$ for three different elements $c, d, e\in [n]\backslash \{x_1, x_2\}$, and $\{x_1, a, b\}, \{x_1, a, y\}\in \mathcal{F}$ for three different elements $a, b, y\in [n]\backslash \{x_1, x_2\}$, then
\begin{align*}
\mathcal{G}\subseteq&\mathcal{R}\cup\left\{\{x_1,c, i\}:  i\in [n]\backslash\{x_1,x_2,c\}\right\}\cup\{\{x_1,e, d\}, \{x_2,b, y\}\}\cup\left\{\{x_2,a, i\}:  i\in [n]\backslash\{x_1,x_2,a\}\right\}.
\end{align*}
So $|\mathcal{G}|\leq 3n-6$.

(ii) If there exist $\{x_2, c, d\}, \{x_2, e, f\}\in \mathcal{F}$ for four different elements $c, d, e, f\in [n]\backslash \{x_1, x_2\}$, and $\{x_1, a, b\}, \{x_1, a, y\}\in \mathcal{F}$ for three different elements $a, b, y\in [n]\backslash \{x_1, x_2\}$, then
\begin{align*}
\mathcal{G}\subseteq&\mathcal{R}\cup\left\{\{x_1,c, i\}:  i\in \{e,f\}\right\}\cup\left\{\{x_1,d, i\}:  i\in \{e,f\}\right\}\cup\left\{\{x_2,a, i\}:  i\in [n]\backslash\{x_1,x_2,a\}\right\}\\
&\cup\{\{x_2,b, y\}\}.
\end{align*}
So $|\mathcal{G}|\leq 2 n$.

(iii) If there exist $\{x_2, c, d\}, \{x_2, e, f\}\in \mathcal{F}$ for four different elements $c, d, e, f\in [n]\backslash \{x_1, x_2\}$, and $\{x_1, a, b\}, \{x_1, y, z\}\in \mathcal{F}$ for four different elements $a, b, y, z\in [n]\backslash \{x_1, x_2\}$, then
\begin{align*}
\mathcal{G}\subseteq&\mathcal{R}\cup\left\{\{x_1,c, i\}:  i\in \{e,f\}\right\}\cup\left\{\{x_1,d, i\}:  i\in \{e,f\}\right\}\cup\left\{\{x_2,a, i\}:  i\in \{y,z\}\right\}\\
&\cup\left\{\{x_2,b, i\}:  i\in \{y,z\}\right\}.
\end{align*}
So $|\mathcal{G}|\leq n+6$.


Returning to case (i), we shall assume that cases (ii) and (iii) are not true. Then  $|\mathcal{F}|\leq n-2+2(n-3)=3n-8$. Therefore, we have
$|\mathcal{F}|+|\mathcal{G}|\leq 6n-14\leq\binom{n-1}{2}+\binom{n-2}{2}+3$.
For case (ii), we shall assume that case (iii) is not true. In addition, we may assume that at least two elements of $\left\{\{x_1,c, i\}:  i\in \{e,f\}\right\}\cup\left\{\{x_1,d, i\}:  i\in \{e,f\}\right\}$ are in $\mathcal{G}$.
Then  $|\mathcal{F}|\leq n-2+n-2+n-3=3n-7$. Therefore, we have
$|\mathcal{F}|+|\mathcal{G}|\leq 5n-7\leq\binom{n-1}{2}+\binom{n-2}{2}+3$.
For case (iii), we may assume that at least two elements of $\left\{\{x_1,c, i\}:  i\in \{e,f\}\right\}\cup\left\{\{x_1,d, i\}:  i\in \{e,f\}\right\}$ are in $\mathcal{G}$, and at least two elements of $\left\{\{x_2,a, i\}:  i\in \{y,z\}\right\}\cup\left\{\{x_2,b, i\}:  i\in \{y,z\}\right\}$ are in $\mathcal{G}$.  Then  $|\mathcal{F}|\leq n-2+n-2+n-2=3n-6$. Therefore, we have
$|\mathcal{F}|+|\mathcal{G}|\leq 4n\leq\binom{n-1}{2}+\binom{n-2}{2}+3$.
This completes the proof.
\end{proof}

\end{document}
